% Shared preamble for the English summaries in this repository.
\documentclass[11pt,a4paper]{article}
\usepackage[margin=2.2cm]{geometry}
\usepackage[T1]{fontenc}
\usepackage[utf8]{inputenc}
\usepackage{lmodern}
\usepackage{amsmath,amssymb}
\usepackage{graphicx}
\usepackage{booktabs}
\usepackage{enumitem}
\usepackage[hidelinks]{hyperref}
\setlist{itemsep=2pt,topsep=4pt}
\setlength{\parindent}{0pt}
\setlength{\parskip}{5pt}
\graphicspath{{figures/}}
\newcommand{\summaryheader}[3]{%
  \begin{center}
    {\Large\bfseries #2}\\[3pt]
    {\large #1}\\[3pt]
    {\small Arshia Goshtasbi \textbullet{} MSc Electrical Engineering (Control), Iran University of Science and Technology \textbullet{} #3}
  \end{center}
  \vspace{2pt}\hrule\vspace{8pt}
  \noindent\textit{English summary of my original Persian report, which is kept as a PDF in this folder.}\par}

\begin{document}
\summaryheader{Multivariable Control, Homework 1}{A Multivariable Industrial Plant: the Boiler--Turbine Unit}{Autumn 1403 (2024)}

\section*{Task}
Introduce an industrial multivariable plant and identify its inputs and outputs.

\section*{What I wrote}
Following Maffezzoni, \emph{Boiler--turbine dynamics in power-plant control} (Control Eng.\ Practice, 1997), I described the boiler--turbine unit of a fossil-fuel power plant as a strongly coupled MIMO system built from three dynamics:
\begin{itemize}
  \item \textbf{Evaporator pressure.} Set by the balance between heat input (mainly from fuel flow) and the energy carried away by steam.
  \item \textbf{Drum level.} Set mainly by the mass balance between feedwater and steam flow, and also disturbed by heat input and pressure swings.
  \item \textbf{Electrical power.} Proportional to high-pressure turbine steam flow, plus a delayed contribution from the reheater.
\end{itemize}
\textbf{Inputs:} fuel flow, air flow, turbine admission (control-valve opening), feedwater flow, recirculated-gas flow, and superheater and reheater spray flows.

\textbf{Outputs:} electrical power, turbine inlet pressure, drum level, turbine inlet temperature, reheater outlet temperature, and air-to-fuel ratio.

Fuel flow changes pressure, which changes drum level, and heat input affects both power and level. Because of this coupling, independent single loops are not enough; the plant needs coordinated multivariable control.
\end{document}
